\section*{Capítulo \ref{ch:b1:faults-folds-fabrics} --- Falhas, dobras e fábricas}

\begin{solution}{exo:b1:faults-folds-fabrics:1}
(a) Uma \omterm{def:b1:faults-folds-fabrics:fault-types}{falha normal}. (b) Uma \omterm{def:b1:faults-folds-fabrics:fault-types}{falha inversa}, quase puramente de \omterm{def:b1:geological-maps:attitude}{mergulho}. (c) Uma \omterm{def:b1:faults-folds-fabrics:fault-types}{falha transcorrente} \omterm{def:b1:faults-folds-fabrics:fault-types}{sinistral}. (d) Uma \omterm{def:b1:faults-folds-fabrics:fault-types}{falha transcorrente} \omterm{def:b1:faults-folds-fabrics:fault-types}{dextral}. (e) Uma \omterm{def:b1:faults-folds-fabrics:fault-types}{falha oblíqua}: o \omterm{def:b1:faults-folds-fabrics:fault-elements}{bloco superior} desceu e deslizou ao longo da \omterm{def:b1:geological-maps:attitude}{direção}, para a esquerda; sua componente transcorrente, $\cos\ang{40} = 0.77$ do \omterm{def:b1:faults-folds-fabrics:fault-elements}{deslizamento líquido}, é maior que sua componente de \omterm{def:b1:geological-maps:attitude}{mergulho}, $0.64$, de modo que é chamada uma falha \omterm{def:b1:faults-folds-fabrics:fault-types}{sinistral} com componente normal. (f) Com um \omterm{def:b1:geological-maps:attitude}{mergulho} de \ang{25}, menor que \ang{45}, a \omterm{def:b1:faults-folds-fabrics:fault-types}{falha inversa} é uma \omterm{def:b1:faults-folds-fabrics:fault-types}{falha de empurrão}.
\end{solution}

\begin{solution}{exo:b1:faults-folds-fabrics:2}
Com $|\sin r| = 1$, o \omterm{def:b1:cross-sections:throw}{rejeito} é $300\sin\ang{55} = \qty{245.7}{m}$ e o salto
$300\cos\ang{55} = \qty{172.1}{m}$.
\end{solution}

\begin{solution}{exo:b1:faults-folds-fabrics:3}
Para uma \omterm{def:b1:faults-folds-fabrics:fold-classes}{dobra vertical} cujos \omterm{def:b1:faults-folds-fabrics:fold-elements}{flancos} mergulham $\delta_1$ e $\delta_2$ em sentidos opostos, $\iota = \ang{180} - \delta_1 - \delta_2$. (a) \ang{110}: aberta.
(b) \ang{35}: fechada. (c) \ang{7}: isoclinal.
\end{solution}

\begin{solution}{exo:b1:faults-folds-fabrics:4}
As falhas normais mergulham $\ang{45} + \ang{35}/2 = \ang{62.5}$, os empurrões
$\ang{45} - \ang{17.5} = \ang{27.5}$. \omterm{def:b1:faults-folds-fabrics:fault-types}{Falhas transcorrentes} conjugadas fazem cada uma
\ang{27.5} com $\sigma_1$, de um lado e do outro dele, de modo que se encontram em \ang{55}.
\end{solution}

\begin{solution}{exo:b1:faults-folds-fabrics:5}
(a) A componente transcorrente é $400\cos\ang{30} = \qty{346.4}{m}$ ao longo da \omterm{def:b1:geological-maps:attitude}{direção}, a componente de \omterm{def:b1:geological-maps:attitude}{mergulho} $400\sin\ang{-30} = \qty{-200}{m}$, isto é, \qty{200}{m} para baixo do \omterm{def:b1:geological-maps:attitude}{mergulho}; o \omterm{def:b1:cross-sections:throw}{rejeito} é $200\sin\ang{50} = \qty{153.2}{m}$ e o salto $200\cos\ang{50} = \qty{128.6}{m}$. (b) A \omterm{def:b1:geological-maps:attitude}{direção} é $000$, logo
$\mathbf s = (0, 1, 0)$; $\mathbf u = (-\cos\ang{50}, 0, \sin\ang{50}) =
(-0.643, 0, 0.766)$. Então
$\overrightarrow{AA'} = 400(0.866\,\mathbf s - 0.5\,\mathbf u) =
(128.6, 346.4, -153.2)$ em metros. Seu \omterm{def:b1:field-lab-methods-1:trend-plunge}{rumo} é $\arctan(128.6/346.4) =
\ang{20.4}$ e seu \omterm{def:b1:field-lab-methods-1:trend-plunge}{caimento} $\arcsin(153.2/400) = \ang{22.5}$, que é também
$\arcsin(\sin\ang{30}\sin\ang{50})$. (c) Uma falha de deslizamento oblíquo: o \omterm{def:b1:faults-folds-fabrics:fault-elements}{rake} jaz
entre $0$ e \ang{-90}, mais perto de $0$: uma falha \omterm{def:b1:faults-folds-fabrics:fault-types}{sinistral} com componente normal.
\end{solution}

\begin{solution}{exo:b1:faults-folds-fabrics:6}
(a) As duas falhas fazem \ang{60} entre si; $\sigma_1$ divide ao meio o ângulo agudo, logo é horizontal ao longo de $050$ (e $230$); $\sigma_3$ é horizontal ao longo de $140$; $\sigma_2$ é vertical, ao longo da linha em que as duas falhas verticais se encontram. Cada falha faz \ang{30} com $\sigma_1$, logo $\ang{45} - \phi/2 =
\ang{30}$ e $\phi = \ang{30}$. (b) As duas falhas cortam o solo em quatro cunhas; as cunhas que contêm $\sigma_1$, em direção a $050$ e a $230$, movem-se uma em direção à outra, e as que contêm $\sigma_3$ afastam-se, em direção a $140$ e $320$. Tome a falha de \omterm{def:b1:geological-maps:attitude}{direção} $020$, entre a cunha em direção a $050$ (a seu leste) e a cunha em direção a $320$ (a seu oeste): em relação à cunha ocidental, a oriental move-se em direção a $230$ e para longe de $320$, ambos com uma componente em direção a $200$ ao longo da falha. Um observador na cunha ocidental, olhando para o leste, vê a oriental mover-se em direção a $200$, para a direita: a falha de \omterm{def:b1:geological-maps:attitude}{direção} $020$ é \omterm{def:b1:faults-folds-fabrics:fault-types}{dextral}. O mesmo argumento, ou a simetria em torno de $\sigma_1$, torna a falha de \omterm{def:b1:geological-maps:attitude}{direção} $080$ \omterm{def:b1:faults-folds-fabrics:fault-types}{sinistral}.
\end{solution}

\begin{solution}{exo:b1:faults-folds-fabrics:7}
As normais ascendentes são
\begin{align*}
  \mathbf n_1 &= (\sin\ang{40}\sin\ang{120}, \sin\ang{40}\cos\ang{120}, \cos\ang{40}) = (0.557, -0.321, 0.766),\\
  \mathbf n_2 &= (\sin\ang{30}\sin\ang{260}, \sin\ang{30}\cos\ang{260}, \cos\ang{30}) = (-0.492, -0.087, 0.866).
\end{align*}
Seu produto vetorial é
$(-0.211, -0.859, -0.207)$, que aponta para baixo: a \omterm{def:b1:faults-folds-fabrics:fold-elements}{linha de charneira} tem \omterm{def:b1:field-lab-methods-1:trend-plunge}{caimento}
$\arcsin 0.207/0.909 = \ang{13.1}$ em direção a $\ang{180} + \arctan(0.211/0.859) =
\ang{193.8}$. $\mathbf n_1\cdot\mathbf n_2 = 0.417$, o ângulo entre as
normais é \ang{65.3}, e $\iota = \ang{114.7}$: uma \omterm{def:b1:faults-folds-fabrics:fold-classes}{dobra aberta}.
\end{solution}

\begin{solution}{exo:b1:faults-folds-fabrics:8}
$\mathbf n_1 = (0.321, 0.383, 0.866)$ para o acamamento e
$\mathbf n_2 = (0.337, -0.925, 0.174)$ para a \omterm{def:b1:faults-folds-fabrics:foliation}{clivagem}. Seu produto vetorial, tomado para baixo, dá uma linha de \omterm{def:b1:field-lab-methods-1:trend-plunge}{caimento} \ang{25.4} em direção a \ang{074.8}. Pela
\cref{prop:b1:faults-folds-fabrics:intersection}, esta é a \omterm{def:b1:geological-maps:attitude}{direção} da \omterm{def:b1:faults-folds-fabrics:fold-elements}{linha de charneira}: a \omterm{def:b1:geological-maps:fold}{dobra} cai cerca de \ang{25} em direção a leste-nordeste, o que qualquer outro \omterm{def:b1:geological-maps:geological-map}{afloramento} de um dos \omterm{def:b1:faults-folds-fabrics:fold-elements}{flancos} deve confirmar.
\end{solution}

\begin{solution}{exo:b1:faults-folds-fabrics:9}
(a) Um observador ao sul da falha, olhando para o norte, vê a parte norte do dique deslocada para leste, para a direita: a falha é \omterm{def:b1:faults-folds-fabrics:fault-types}{dextral}. (b) O dique é um plano vertical perpendicular à falha, logo seu deslocamento ao longo da \omterm{def:b1:geological-maps:attitude}{direção} é a componente transcorrente: \qty{350}{m}, que é o deslizamento inteiro se as \omterm{def:b1:faults-folds-fabrics:fault-rocks}{estrias} forem horizontais. (c) O deslizamento então faz \ang{20} com a \omterm{def:b1:geological-maps:attitude}{direção} no plano vertical da falha: sua componente horizontal continua sendo \qty{350}{m}, logo o \omterm{def:b1:faults-folds-fabrics:fault-elements}{deslizamento líquido} é $350/\cos\ang{20} = \qty{372.5}{m}$ e sua componente vertical $350\tan\ang{20} = \qty{127.4}{m}$. O bloco norte moveu-se para leste, e as \omterm{def:b1:faults-folds-fabrics:fault-rocks}{estrias} descem para o leste, logo ele também desceu, em \qty{127.4}{m}. Um dique vertical não pode mostrar esse movimento vertical, já que deslizá-lo ao longo de si mesmo deixa seu traço inalterado.
\end{solution}

\begin{solution}{exo:b1:faults-folds-fabrics:10}
(a) A normal à falha faz $\ang{90} - \theta$ com $\sigma_1$, logo, pela
\cref{prop:b1:stress-strain-rheology:mohr},
$\sigma_n' = \sigma_1'\sin^2\theta + \sigma_3'\cos^2\theta$ e
$\tau = (\sigma_1' - \sigma_3')\sin\theta\cos\theta$. O deslizamento começa quando
$\tau = \mu\sigma_n'$; dividindo por $\sigma_3'\cos^2\theta$ e escrevendo
$R = \sigma_1'/\sigma_3'$, $(R - 1)\tan\theta = \mu(R\tan^2\theta + 1)$, logo
$R\tan\theta(1 - \mu\tan\theta) = \tan\theta + \mu$, que é a fórmula.
(b) Com $\sigma_1$ vertical, $\theta$ é \ang{90} menos o \omterm{def:b1:geological-maps:attitude}{mergulho}:
$\theta = \ang{30}$ dá $R = 3.12$; $\theta = \ang{45}$, $R = 4.00$;
$\theta = \ang{55}$, $R = 9.92$. (c) O denominador se anula quando
$\tan\theta = 1/\mu$, $\theta = \ang{59.0}$; além disso, $R$ seria negativo:
nenhuma razão de tensões efetivas positivas pode fazer a falha deslizar, por maior que
$\sigma_1'$ se torne. Falhas normais que mergulham menos que \ang{31.0} não podem deslizar
sob um $\sigma_1$ vertical com $\mu = 0.6$. (d) O plano mais bem orientado está a
$\ang{45} - \phi/2$ de $\sigma_1$. Para $\mu = 0.6$, $\phi = \ang{31.0}$:
falhas normais mergulham \ang{60.5}, empurrões \ang{29.5}. Para $\mu = 0.85$,
$\phi = \ang{40.4}$: \ang{65.2} e \ang{24.8}.
\end{solution}

\begin{solution}{exo:b1:faults-folds-fabrics:11}
Olhando para o norte, o leste está à direita, e planos que mergulham para oeste sobem em direção à direita.
(a) Em B o acamamento mergulha no mesmo sentido que a \omterm{def:b1:faults-folds-fabrics:foliation}{clivagem}, e mais íngreme: pela
\cref{prop:b1:faults-folds-fabrics:cleavage-bedding}(a), B está em um flanco invertido; A e C, onde o acamamento é menos íngreme, estão na posição normal. (b) A \omterm{def:b1:faults-folds-fabrics:foliation}{clivagem} sobe a \ang{60}; em A e C o acamamento sobe a
\ang{40}, e a rotação aguda de um para o outro é horária, por
\ang{20}, movendo o topo da \omterm{def:b1:faults-folds-fabrics:foliation}{clivagem} para leste; em B o acamamento sobe a
\ang{80} e a rotação é anti-horária, por \ang{20}, movendo-o para o oeste. (c) O sentido muda entre A e B e de novo entre B e C: há uma \omterm{def:b1:faults-folds-fabrics:fold-elements}{charneira} em cada intervalo. Em A, que está na posição normal, o topo da \omterm{def:b1:faults-folds-fabrics:foliation}{clivagem} se move para leste, em direção ao \omterm{def:b1:geological-maps:fold}{anticlinal} no topo do flanco: a \omterm{def:b1:faults-folds-fabrics:fold-elements}{charneira} entre A e B é um \omterm{def:b1:geological-maps:fold}{anticlinal}, e a que fica entre B e C, um \omterm{def:b1:geological-maps:fold}{sinclinal}.
(d) Na notação da prova, os \omterm{def:b1:faults-folds-fabrics:fold-elements}{flancos} na posição normal mergulham $\gamma - \iota/2 = \ang{40}$ e os \omterm{def:b1:faults-folds-fabrics:fold-elements}{flancos} invertidos $\gamma + \iota/2 =
\ang{80}$, logo $\gamma = \ang{60}$, como diz a \omterm{def:b1:faults-folds-fabrics:foliation}{clivagem}, e $\iota =
\ang{40}$: \omterm{def:b1:faults-folds-fabrics:fold-classes}{dobras fechadas}. Os planos axiais têm \omterm{def:b1:geological-maps:attitude}{direção} norte e mergulham \ang{60} para oeste, \omterm{def:b1:geological-maps:attitude}{atitude} $270/60$. Os \omterm{def:b1:faults-folds-fabrics:fold-elements}{flancos} curtos passaram da vertical com seus topos para leste: as \omterm{def:b1:geological-maps:fold}{dobras} vergem para leste, como na figura do capítulo.
\end{solution}

\begin{solution}{exo:b1:faults-folds-fabrics:12}
(a) Numa \omterm{def:b1:geological-maps:fold}{dobra} em chevron simétrica, cada flanco faz o ângulo $\iota/2$ com o \omterm{def:b1:faults-folds-fabrics:fold-elements}{plano axial}, que é perpendicular à \omterm{def:b1:geological-maps:attitude}{direção} de encurtamento; um flanco de comprimento $l$ ocupa portanto $l\sin(\iota/2)$ através das \omterm{def:b1:geological-maps:fold}{dobras}, e a camada inteira $L_0\sin(\iota/2)$.
(b) $\sin\ang{30} = 0.5$: um encurtamento de
\qty{50}{\percent}. (c) $\sin(\iota/2) = 0.7$ dá $\iota = \ang{88.9}$.
(d) $\iota = \ang{156}$ e $\sin\ang{78} = 0.978$: um encurtamento de
\qty{2.2}{\percent}, que é também $1 - \cos\ang{12}$.
\end{solution}

\begin{solution}{pb:b1:faults-folds-fabrics:1}
\textbf{1.} $\mathbf n = (\sin\delta\sin\alpha, \sin\delta\cos\alpha,
\cos\delta)$ dá
\[
  \mathbf n_1 = (0.20855, 0.03415, 0.97742), \qquad
  \mathbf n_2 = (-0.20855, 0.03415, 0.97742).
\]

\textbf{2.} $\mathbf n_1\cdot\mathbf n_2 = -0.04349 + 0.00117 + 0.95535 =
0.91302$, e o ângulo entre as normais é \ang{24.1}.

\textbf{3.} Pela \cref{prop:b1:faults-folds-fabrics:hinge},
$\iota = \ang{180} - \ang{24.1} = \ang{155.9}$.

\textbf{4.} Entre \ang{170} e \ang{90}: uma \omterm{def:b1:faults-folds-fabrics:fold-classes}{dobra aberta}.

\textbf{5.} Os \omterm{def:b1:faults-folds-fabrics:fold-elements}{flancos} mergulham em direção a $080.7$ e $279.3$, e não exatamente para leste e oeste, porque a \omterm{def:b1:geological-maps:fold}{dobra} tem \omterm{def:b1:field-lab-methods-1:trend-plunge}{caimento}: suas linhas de maior declive não estão no perfil, que mostra apenas um \omterm{def:b1:cross-sections:apparent-dip}{mergulho aparente}. Em uma seção leste--oeste, pela
\cref{thm:b1:cross-sections:apparent-dip}, $\tan\beta = \tan\ang{12.2}
\cos\ang{9.3} = 0.2134$, logo $\beta = \ang{12.0}$, e $\ang{180} -
2\times\ang{12.0} = \ang{155.9}$ (o perfil está inclinado \ang{2} dessa seção, o que não altera nada com esta precisão).

\textbf{6.} Seus \omterm{def:b1:faults-folds-fabrics:fold-elements}{flancos} são paralelos aos do \omterm{def:b1:geological-maps:fold}{anticlinal}, logo seu \omterm{def:b1:faults-folds-fabrics:fold-elements}{ângulo entre flancos} também é \ang{155.9}.

\textbf{7.} $\mathbf n_1\times\mathbf n_2 = (0.03415\times0.97742 -
0.97742\times0.03415,\; 0.97742\times(-0.20855) - 0.20855\times0.97742,\;
0.20855\times0.03415 + 0.03415\times0.20855) = (0, -0.40768, 0.01424)$.
Tomado para baixo, $(0, 0.40768, -0.01424)$: a \omterm{def:b1:faults-folds-fabrics:fold-elements}{linha de charneira} tem \omterm{def:b1:field-lab-methods-1:trend-plunge}{rumo} $000$ e
tem \omterm{def:b1:field-lab-methods-1:trend-plunge}{caimento} $\arctan(0.01424/0.40768) = \ang{2.0}$ para o norte.

\textbf{8.} O \omterm{def:b1:faults-folds-fabrics:fold-elements}{plano axial} contém a \omterm{def:b1:faults-folds-fabrics:fold-elements}{linha de charneira} e biparte os dois \omterm{def:b1:faults-folds-fabrics:fold-elements}{flancos};
sua normal está ao longo de $\mathbf n_1 - \mathbf n_2 = (0.41710, 0, 0)$, horizontal
e leste--oeste. Ele é vertical e tem \omterm{def:b1:geological-maps:attitude}{direção} norte--sul: $090/90$. O \omterm{def:b1:geological-maps:fold}{anticlinal} é vertical.

\textbf{9.} Uma \omterm{def:b1:faults-folds-fabrics:fold-classes}{dobra aberta} e vertical, com \omterm{def:b1:field-lab-methods-1:trend-plunge}{caimento} de \ang{2.0} para o norte.

\textbf{10.} Pela \cref{prop:b1:geological-maps:folds}, os \omterm{def:b1:geological-maps:geological-map}{afloramentos} de um \omterm{def:b1:geological-maps:fold}{anticlinal} fecham-se na \omterm{def:b1:geological-maps:attitude}{direção} do \omterm{def:b1:field-lab-methods-1:trend-plunge}{caimento}: para o norte. Seu núcleo contém as rochas mais antigas: o arenito na base da cobertura e, onde o rio cortou mais fundo, o embasamento de xisto e granito sob ela.

\textbf{11.} Plano-axial, a \omterm{def:b1:faults-folds-fabrics:foliation}{clivagem} é vertical e tem \omterm{def:b1:geological-maps:attitude}{direção} norte--sul, $090/90$. Pela \cref{prop:b1:faults-folds-fabrics:intersection}, sua \omterm{def:b1:faults-folds-fabrics:lineation}{lineação de interseção} com o acamamento é paralela à \omterm{def:b1:faults-folds-fabrics:fold-elements}{linha de charneira}: \omterm{def:b1:field-lab-methods-1:trend-plunge}{rumo} $000$, \omterm{def:b1:field-lab-methods-1:trend-plunge}{caimento} \ang{2.0}. Diretamente: o plano $080.7/12.2$ encontra o plano vertical de \omterm{def:b1:geological-maps:attitude}{direção} norte ao longo de uma linha de \omterm{def:b1:field-lab-methods-1:trend-plunge}{caimento} $\arctan(\tan\ang{12.2}
\cos\ang{80.7}) = \ang{2.0}$ em direção a $000$.

\textbf{12.} Olhando para o norte, o flanco leste sobe em direção à esquerda com \ang{12.0}, e a \omterm{def:b1:faults-folds-fabrics:foliation}{clivagem} é vertical: a rotação aguda da \omterm{def:b1:faults-folds-fabrics:foliation}{clivagem} para o acamamento, por \ang{78.0}, é anti-horária e move o topo da \omterm{def:b1:faults-folds-fabrics:foliation}{clivagem} para oeste. O flanco está na posição normal, logo, pela \cref{prop:b1:faults-folds-fabrics:cleavage-bedding}(b), a \omterm{def:b1:faults-folds-fabrics:fold-elements}{charneira} do \omterm{def:b1:geological-maps:fold}{anticlinal} fica a oeste, em $x = \qty{3.5}{km}$, como mostra o mapa.

\textbf{13.} A \omterm{def:b1:geological-maps:attitude}{direção} é $120 - 90 = 030$; o plano mergulha em direção a $120$, logo o \omterm{def:b1:faults-folds-fabrics:fault-elements}{bloco superior} jaz do lado leste-sudeste do traço.

\textbf{14.} A falha é normal: $\sigma_1$ era vertical.

\textbf{15.} $\sigma_2$ é horizontal ao longo da \omterm{def:b1:geological-maps:attitude}{direção}, $030$--$210$; $\sigma_3$ horizontal através dela, $120$--$300$.

\textbf{16.} $\ang{45} + \phi/2 = \ang{60}$ dá $\phi = \ang{30}$, e
$\mu = \tan\ang{30} = 0.577$.

\textbf{17.} O conjunto conjugado tem também \omterm{def:b1:geological-maps:attitude}{direção} $030$ e mergulha \ang{60} para o outro lado: $300/60$. Os dois conjuntos se encontram em $\ang{180} - 2\times\ang{60} = \ang{60}$, com a $\sigma_1$ vertical dividindo ao meio o ângulo agudo.

\textbf{18.} Com $\mu = 0.85$, $\phi = \ang{40.4}$ e as falhas normais mais bem orientadas mergulham \ang{65.2}; a falha de Varden, a \ang{60}, está a \ang{5.2} de distância. Pelo \cref{exo:b1:faults-folds-fabrics:10}, ela desliza na razão $\sigma_1'/\sigma_3' = (1 + 0.85\cot\ang{30})/(1 - 0.85\tan\ang{30}) = 4.85$, contra $4.68$ para o plano mais bem orientado: ela está quase, mas não exatamente, na melhor orientação, e deslizaria de novo com uma razão de tensões apenas \qty{4}{\percent} maior.

\textbf{19.} O deslizamento é para baixo do \omterm{def:b1:geological-maps:attitude}{mergulho}: $r = \ang{-90}$; o \omterm{def:b1:faults-folds-fabrics:fault-elements}{vetor de deslizamento} tem \omterm{def:b1:field-lab-methods-1:trend-plunge}{rumo}
$120$ e \omterm{def:b1:field-lab-methods-1:trend-plunge}{caimento} \ang{60}.

\textbf{20.} $(\cos\ang{60}\sin\ang{120}, \cos\ang{60}\cos\ang{120},
-\sin\ang{60}) = (0.4330, -0.2500, -0.8660)$.

\textbf{21.} $e = 200/\tan\ang{60} = \qty{115.5}{m}$.

\textbf{22.} Leste: $115.47\sin\ang{120} = \qty{100.0}{m}$; norte:
$115.47\cos\ang{120} = \qty{-57.7}{m}$, isto é, \qty{57.7}{m} para o sul.

\textbf{23.} $L = 200/\sin\ang{60} = \qty{230.9}{m}$, e
$200^2 + 115.47^2 = 40\,000 + 13\,333 = 53\,333 = 230.94^2$. O \omterm{def:b1:cross-sections:throw}{rejeito} e o salto são as componentes vertical e horizontal do deslizamento, que jaz no plano vertical do \omterm{def:b1:geological-maps:attitude}{mergulho}: são os catetos de um triângulo retângulo cuja hipotenusa é o deslizamento.

\textbf{24.} Com aproximação de um metro, o \omterm{def:b1:faults-folds-fabrics:fault-elements}{deslizamento líquido} na falha de Varden é
\textbf{\qty[text-series-to-math=true]{231}{m}}.
\end{solution}
